\begin{figure}[ht]
\centering
\hspace{-0.5cm}
\begin{tikzpicture}
\begin{axis}[
width=0.98\columnwidth,
xmin=0.0, xmax=1.95,
ymin=45.65, ymax=47.40,
ylabel={Performance},
xlabel={Time per sample (ms)},
legend columns=1,
legend cell align=left,
legend style={
    anchor=south,
    at={(0.84, 0.05)},
},
]
\addplot[only marks, mark=*, mark options={draw=mhcolor, fill=mhcolor, scale=2}] coordinates {
    (0.372, 45.95)
    (1.514, 47.206)
};
\node at (axis cs:0.372,45.95) [anchor= north west, color=darkgrey] {\small{MHA-Large}};            
\node at (axis cs:1.514,47.206) [anchor= north east, color=darkgrey] {\small{MHA-XXL}};            
\addplot[only marks, mark=*, mark options={draw=mqcolor, fill=mqcolor, scale=2}] coordinates {
    (0.239, 46.571)
};
\node at (axis cs:0.239,46.548) [anchor= north west, color=darkgrey] {\small{MQA-XXL}};            
\addplot[only marks, mark=*, mark options={draw=gqcolor, fill=gqcolor, scale=2}] coordinates {
    (0.275, 47.136)
};
\node at (axis cs:0.275,47.136) [anchor= north west, color=darkgrey] {\small{GQA-XXL}};            
\end{axis}
\end{tikzpicture}
    \caption{{\bf 与 \mh 相比，经过继续训练的 \mq 实现了有利的权衡，其质量高于且速度快于 \mh-Large，而 \gmq 在获得相近速度提升的同时实现了更好的性能，质量与 \mh-XXL 相当。} 对采用 multi-head 注意力的 T5-Large 和 T5-XXL，以及经过 5\% 继续训练、采用 \mq 和 \gmq-8 注意力的 T5-XXL，展示所有任务平均性能随每个样本平均推理时间的变化。}
    \label{fig:perf_vs_time}
\end{figure}
